\documentclass[11pt]{article}
\usepackage[margin=1in]{geometry}
\usepackage{parskip}
\usepackage[colorlinks=true,urlcolor=blue]{hyperref}
\begin{document}
\pagestyle{empty}

\noindent Krishna Harish\\
Lawrence E. Elkins High School, Missouri City, Texas, USA\\
krishnaharish2009@gmail.com

\vspace{1em}
\noindent\textbf{Manuscript ID: ao-2026-07482d (revised submission)}

\vspace{1em}
\noindent Dear Prof.\ Skorb,

Please find enclosed the revised version of ``Objective-Dependent Active Learning and
Calibrated Uncertainty for Sample-Efficient Discovery of Rare-Earth Extractants: A
Benchmark on 1{,}202 Experimental Distribution Coefficients with Prospective
Validation'' (ao-2026-07482d), together with a point-by-point response to the reviewer,
a change-highlighted copy, and a clean copy.

We are grateful to the reviewer for a detailed and constructive report. We have
addressed every point with new experiments on the real data; no claim is made without a
computed result. The revision is substantial: the manuscript now reports five findings
(adding a leakage-controlled generalization analysis and a definition-dependent
discovery analysis), four new tables, and five new figures, all regenerated by released
scripts. The principal additions, each responding to a specific reviewer request, are:

\begin{enumerate}
\item \textbf{Leakage-controlled evaluation} (Reviewer 1.1): recovering the ligand,
source, and lanthanide grouping structure from the released features, we show that
random row-level cross-validation ($R^2=0.88$) is strongly optimistic: ligand-, source-,
and structural-family-disjoint splits collapse the estimate to at or below the mean. The
in-distribution $R^2$ is now explicitly reframed as not a deployability estimate.
\item \textbf{Ligand-batch acquisition} (Reviewer 1.2): the realistic laboratory
decision unit (select a ligand, measure all its conditions); the objective-dependent
dichotomy persists.
\item \textbf{Definition-resolved discovery} (Reviewer 1.3): recall reported for rows,
ligand--metal pairs, and unique ligands, with enrichment factor, regret, and hit-rate;
the strategy ranking inverts across definitions.
\item \textbf{Rigorous subgroup calibration} (Reviewer 1.4): ligand-disjoint
train/calibration/test; coverage by logD region and lanthanide class; Mondrian
group-conditional conformal; and a demonstration that non-monotone calibration reorders
the uncertainty ranking.
\item \textbf{Broader model and uncertainty-estimator baselines} (Reviewer 1.5): random
forest, quantile regression forest, Gaussian process, gradient boosting, and an
MC-dropout Bayesian neural network, plus a selectivity (separation-factor) analysis
showing a high distribution coefficient alone is not sufficient.
\end{enumerate}

We have also added figures and tables throughout (Reviewer 1.6, 1.7) and expanded the
literature with 17 verified, up-to-date references (Reviewer 1.8).

All code, recovered group labels, and split definitions are released at
\url{https://github.com/krishnatheaverage/logd-active-learning-benchmark} and archived on
Zenodo; every reported number is reproducible on a single laptop. The manuscript is
original and not under consideration elsewhere; use of an AI tool is disclosed within;
the author has no competing interests to declare, and there are no funding sources to
report.

\vspace{1em}
\noindent Thank you for the opportunity to revise.

\vspace{1em}
\noindent Sincerely,\\
Krishna Harish

\end{document}
